\documentclass[../../main.tex]{subfiles}
\begin{document}

{\bf [解答]}

$C$上の点$(X(\theta),Y(\theta))$とすると，題意より
\begin{align}
  \begin{cases}
    X(\theta)=(2+\cos\theta)\cos\theta \\
    Y(\theta)=(2+\cos\theta)\sin\theta
  \end{cases} \label{eq:1}
\end{align}
である．$\theta$による一階微分は
\begin{align}
  \begin{cases}
    X'(\theta)=-2\cos\theta\sin\theta-2\sin\theta=-2\sin\theta(\cos\theta+1) \\
    Y'(\theta)=2\cos\theta-\sin^2\theta+\cos^2\theta = 2\cos^2\theta+2\cos\theta-1
  \end{cases} \label{eq:2}
\end{align}
である．従って$0\leqq\theta\leqq\pi$では$X'\leqq0$かつ$Y\geqq0$である．
また，端点では$\theta=0$で$P(3,0)$，$\theta=\pi$で$P'(-1,0)$であり，増減表および$C$の概形は以下のようになる．

\begin{table}[htb]
  \centering
  \caption{曲線 $C$ の増減表（$0\leqq\theta\leqq\pi$）}
  \label{tab:curve_C}
  \begin{tabular}{c|ccccc}
    $\theta$ & $0$ & $\cdots$ & $\alpha$ & $\cdots$ & $\pi$ \\
    \hline
    $X'$ & $0$ & $-$ & $-$ & $-$ & $0$ \\
    \hline
    $X$ & $3$ & $\searrow$ & $\dfrac{\sqrt{3}}{2}$ & $\searrow$ & $-1$ \\
    \hline
    $Y'$ & $3$ & $+$ & $0$ & $-$ & $-1$ \\
    \hline
    $Y$ & $0$ & $\nearrow$ & 最大 & $\searrow$ & $0$
  \end{tabular}
  \\[2pt]
  \footnotesize（ただし $\alpha$ は $\cos\alpha = \dfrac{\sqrt{3}-1}{2}$（$0<\alpha<\pi$）を満たす角）
\end{table}

\begin{figure}[htb]
  \centering
  \begin{tikzpicture}[x=1.3cm, y=1.3cm, >=stealth, every node/.style={font=\small}]

    % 座標軸
    \draw[->, thick, gray!80!black] (-1.8, 0) -- (4.8, 0) node[right] {$x$};
    \draw[->, thick, gray!80!black] (0, -0.8) -- (0, 3.2) node[above] {$y$};
    \node[below left=2pt, font=\footnotesize] at (0, 0) {$O$};

    % 回転体の下半分の輪郭（薄い破線で立体の奥行きを表現）
    \draw[dashed, line width=0.8pt, blue!40!gray, domain=0:180, samples=80]
    plot ({ (2 + cos(\x))*cos(\x) }, { -(2 + cos(\x))*sin(\x) });

    % 回転断面の楕円（x = sqrt(3)/2 付近）
    \draw[dotted, gray!50, line width=0.6pt]
    (0.866, 0) ellipse (0.35cm and 2.2cm);

    % C と x 軸で囲まれた領域の塗り
    \fill[blue!15, opacity=0.6]
    plot[domain=0:180, samples=80] ({ (2 + cos(\x))*cos(\x) }, { (2 + cos(\x))*sin(\x) })
    -- (-1, 0) -- (3, 0) -- cycle;

    % 曲線 C (0 <= theta <= pi)
    \draw[line width=1.4pt, blue!85!black, domain=0:180, samples=100]
    plot ({ (2 + cos(\x))*cos(\x) }, { (2 + cos(\x))*sin(\x) });
    \node[above right=1pt, blue!85!black, font=\normalsize\bfseries] at (2.4, 1.4) {$C$};

    % 曲線上の移動方向矢印 (theta: 0 -> pi)
    \draw[->, line width=1.0pt, teal!80!black]
    ({(2+cos(40))*cos(40)}, {(2+cos(40))*sin(40)}) --
    ({(2+cos(50))*cos(50)}, {(2+cos(50))*sin(50)});
    \draw[->, line width=1.0pt, teal!80!black]
    ({(2+cos(115))*cos(115)}, {(2+cos(115))*sin(115)}) --
    ({(2+cos(125))*cos(125)}, {(2+cos(125))*sin(125)});
    \node[teal!80!black, font=\scriptsize, above=2pt]
    at ({(2+cos(48))*cos(48)}, {(2+cos(48))*sin(48)}) {$\theta$ 増加方向};

    % 特徴点
    \coordinate (P) at (3, 0);
    \coordinate (Pprime) at (-1, 0);
    \coordinate (Top) at (0.866, 2.20);
    \coordinate (YaxisPt) at (0, 2);

    % 補助線
    \draw[dashed, gray!60, line width=0.7pt] (Top) -- (0.866, 0);
    \node[circle, fill, inner sep=1.0pt] at (0.866, 0) {};
    \node[below=2pt, font=\scriptsize] at (0.866, 0) {$\frac{\sqrt{3}}{2}$};

    % y 軸との交点 (0, 2)
    \node[circle, fill=teal!80!black, inner sep=1.4pt] at (YaxisPt) {};
    \node[left=2pt, font=\scriptsize, teal!80!black] at (YaxisPt) {$2$};

    % 点ノード
    \node[circle, fill=blue!85!black, inner sep=1.5pt] at (P) {};
    \node[above right=2pt, font=\footnotesize] at (P) {$P(3,0)$};
    \node[below=3pt, font=\scriptsize, gray!80!black] at (P) {($\theta=0$)};

    \node[circle, fill=blue!85!black, inner sep=1.5pt] at (Pprime) {};
    \node[above left=2pt, font=\footnotesize] at (Pprime) {$P'(-1,0)$};
    \node[below=3pt, font=\scriptsize, gray!80!black] at (Pprime) {($\theta=\pi$)};

    % 極大点
    \node[circle, fill=red!85!black, inner sep=1.5pt] at (Top) {};
    \node[above=3pt, font=\footnotesize, red!85!black] at (Top) {極大点 $\left(\frac{\sqrt{3}}{2}, Y_{\max}\right)$};

    % x 軸周りの回転矢印
    \draw[->, line width=1.1pt, orange!90!black]
    (3.8, 0.45) arc[start angle=60, end angle=300, x radius=0.25cm, y radius=0.45cm];
    \node[below=2pt, font=\scriptsize, orange!90!black] at (3.8, -0.45) {$x$軸回転};

  \end{tikzpicture}
  \caption{曲線 $C: r=2+\cos\theta$（$0\leqq\theta\leqq\pi$）の概形と回転体の領域}
  \label{fig:curve_C}
\end{figure}

従って求める体積$V$は
\begin{align}
  \begin{split}
    V
    &=\int_{-1}^3\pi Y^2\,dX\\
    &=\int_{\pi}^0\pi\{(2+\cos\theta)\sin\theta\}^2\dfrac{dX}{d\theta}\,d\theta\\
    &=\pi\int_0^\pi(2+\cos\theta)^2\sin^2\theta\,2\sin\theta(\cos\theta+1)\,d\theta\\
    &=2\pi\int_0^\pi(2+\cos\theta)^2(1-\cos^2\theta)(\cos\theta+1)\sin\theta\,d\theta\\
    &=2\pi\int_0^\pi(-\cos^5\theta-5\cos^4\theta-7\cos^3\theta+\cos^2\theta+8\cos\theta+4)\sin\theta\,d\theta\\
    &=-2\pi\left[-\dfrac{1}{6}\cos^6\theta-\cos^5\theta-\dfrac{7}{4}\cos^4\theta
    +\dfrac{1}{3}\cos^3\theta+4\cos^2\theta+4\cos\theta\right]_0^{\pi}\\
    &=-2\pi\cdot2\left(1-\dfrac{1}{3}-4\right)\\
    &=\dfrac{40}{3}\pi
  \end{split}
\end{align}
である．

\end{document}
